\noindent340\hfill\textbf{BEMERKUNGEN.}\hfill\mbox{}
\begin{center}\textbf{BEMERKUNGEN.}\end{center}

\par\noindent\hspace*{\parindent}Die vorstehende unvollendete Ausgleichung ist der einzige Eintrag in einem Notizbuch mit dem Titel »Rechnungen in Beziehung auf die trigonometrischen Messungen«. Sie ist bemerkenswert, weil hier zum ersten Mal die Polygongleichungen für ein Kranzsystem aufgestellt werden. Die 3 Polygongleichungen bestehen aus der Winkelgleichung für das innere Polygon des Kranzes und den beiden Bedingungsgleichungen, die sicherstellen, dass die Koordinaten eines Punktes wieder dieselben Werte annehmen, wenn man auf einem aus Seiten des Kranzsystems gebildeten Linienzug zu ihm zurückkehrt. Diese beiden Bedingungen werden durch Gleichung 2) auf S.~337 dargestellt, die sich wie folgt herleiten lässt.

\par\noindent\hspace*{\parindent}Berechnet man vom Punkt 1 aus entlang des Linienzuges 1. 21. 3. 4. 5. 6. 9. 10. 11. 12. 15. 16. 18. 19. 20 die rechtwinkligen Koordinaten der Punkte 21 und 1, so erhält man vor der Ausgleichung im Allgemeinen Werte, die von den Ausgangswerten für 21 und 1 abweichen und den Punkten \(21^*\) und \(1^*\) entsprechen mögen. Bezeichnet man die Lage des Punktes \(k\) mit \(z_k=x_k+\mathrm{i}y_k\), so gilt
\[
\begin{aligned}
z_{1^*}-z_1={}&(z_{21}-z_1)+(z_3-z_{21})+(z_4-z_3)+\ .\ .\ .\\
&+(z_{20}-z_{19})+(z_{21^*}-z_{20})+(z_{1^*}-z_{21^*}).
\end{aligned}
\]
\par\noindent\hspace*{\parindent}Um \(z_{1^*}\) in \(z_1\) zu überführen, werden die Seiten und Winkel des Linienzuges differenziell geändert. Dabei wird der Punkt 1 als fest angenommen.
\[
\delta z_{1^*}=\delta(z_{21}-z_1)+\delta(z_3-z_{21})+\ .\ .\ .+\delta(z_{1^*}-z_{21^*}).
\]
\par\noindent\hspace*{\parindent}Ist \(\mathrm{s}_{\lambda\mu}\) die Länge und \(\tau_{\lambda.\mu}\) das Azimut der Seite \(\lambda.\mu\), so gilt
\[
z_\mu-z_\lambda=\mathrm{s}_{\lambda.\mu}\mathrm{e}^{\mathrm{i}\tau_{\lambda.\mu}}
\]
\noindent und
\[
\delta(z_\mu-z_\lambda)=(z_\mu-z_\lambda)(\delta\log\mathrm{s}_{\lambda.\mu}+\mathrm{i}\,.\,\delta\tau_{\lambda.\mu}),
\]
\noindent sodass
\[
1)\quad\delta z_{1^*}=\Sigma(z_\mu-z_\lambda)\delta\log\mathrm{s}_{\lambda.\mu}+\mathrm{i}\Sigma(z_\mu-z_\lambda)\delta\tau_{\lambda.\mu},
\]
\noindent wobei sich die Summen auf die Punkte des Linienzuges beziehen. Der Logarithmus ist hier wie auch im Folgenden der hyperbolische Logarithmus.

\par\noindent\hspace*{\parindent}Bezeichnet man den Winkel zwischen den aufeinanderfolgenden Seiten \(\lambda.\mu\) und \(\mu.\nu\) vorübergehend mit \(w_\mu\), so gilt
\[
\tau_{\mu.\nu}=\tau_{\lambda.\mu}+w_\mu\pm180^0
\]
\noindent und daher
\[
\delta\tau_{\mu.\nu}=\delta\tau_{\lambda.\mu}+\delta w_\mu;
\]
\noindent Drückt man mithilfe dieser Gleichung nacheinander alle \(\delta\tau\) durch \(\delta\tau_{1.21}\) aus, so erhält man für den Linienzug
\[
\begin{aligned}
\Sigma(z_\mu-z_\lambda)\delta\tau_{\lambda.\mu}={}&(z_{1^*}-z_1)\delta\tau_{1.21}+(z_{1^*}-z_{21})\delta w_{21}\\
&+(z_{1^*}-z_3)\delta w_3+\ .\ .\ .+(z_{1^*}-z_{20})\delta w_{20}\\
&+(z_{1^*}-z_{21^*})\delta w_{21^*}.
\end{aligned}
\]
\par\noindent\hspace*{\parindent}\(z_{1^*}-z_1\) kann gegenüber den anderen Differenzen \(z_{1^*}-z_\lambda\) als kleine Größe erster Ordnung angesehen
\par\clearpage
\noindent\textbf{DREIECKSKRANZ UM OLDENBURG.}\hfill341\par\medskip
\enlargethispage{1.5\baselineskip}
\begingroup\small\setlength{\jot}{1pt}\setlength{\abovedisplayskip}{4pt}\setlength{\belowdisplayskip}{4pt}
\setlength{\abovedisplayshortskip}{2pt}\setlength{\belowdisplayshortskip}{3pt}
\noindent werden. Daher darf man das erste Glied gegenüber den übrigen vernachlässigen und in den Koeffizienten \(z_{1^*}\) und \(z_{21^*}\) durch \(z_1\) und \(z_{21}\) ersetzen. Die Abbildung auf S.~332 liefert außerdem
\[
\begin{gathered}
w_{21}=a^0+b^0,\quad w_3=360^0-c^0,\quad w_4=d^0,\\
w_5=360^0-e^0,\quad w_6=f^0+g^0+h^0,\ \text{usw.,}\\
w_{20}=360^0-u^0,\qquad w_{21}=v^0
\end{gathered}
\]
\noindent sodass sich ergibt
\[
\begin{aligned}
&2)\quad\Sigma(z_\mu-z_\lambda)\delta\tau_{\lambda.\mu}=(z_1-z_{21})(\delta a^0+\delta b^0+\delta v^0)\\
&-(z_1-z_3)\delta c^0+(z_1-z_4)\delta d^0-(z_1-z_5)\delta e^0\\
&+(z_1-z_6)(\delta f^0+\delta g^0+\delta h^0)\!+\ .\ .\ .-(z_1-z_{20})\delta u^0.
\end{aligned}
\]
\par\noindent\hspace*{\parindent}Auch die erste Summe in Gleichung (1) lässt sich umformen.

\par\noindent\hspace*{\parindent}Aus
\[
\begin{aligned}
\frac{\mathrm{s}_{21.3}}{\mathrm{s}_{1.21}}&=\frac{\sin b'}{\sin b''}\,.\,\frac{\sin a'}{\sin a''}\\[7pt]
\frac{\mathrm{s}_{3.4}}{\mathrm{s}_{21.3}}&=\frac{\sin c'}{\sin c''}\\[5pt]
&\text{usw.}
\end{aligned}
\]
\noindent folgt
\[
\begin{aligned}
\delta\log\mathrm{s}_{21.3}={}&\delta\log\mathrm{s}_{1.21}+\operatorname{cotang}a'\,.\,\delta a'-\operatorname{cotang}a''\,.\,\delta a''\\
&+\operatorname{cotang}b'\,.\,\delta b'-\operatorname{cotang}b''\,.\,\delta b''\\
\delta\log\mathrm{s}_{3.4}={}&\delta\log\mathrm{s}_{21.3}+\operatorname{cotang}c'\,.\,\delta c'-\operatorname{cotang}c''\,.\,\delta c''\\
&.\quad.\quad.\quad.\quad.\quad.\quad.\quad.\quad.\quad.\quad.\quad.\quad.\quad.\quad.\quad.\quad.\quad.\quad.\quad.\quad.\\
&.\quad.\quad.\quad.\quad.\quad.\quad.\quad.\quad.\quad.\quad.\quad.\quad.\quad.\quad.\quad.\quad.\quad.\quad.\quad.\quad.\\
\delta\log\mathrm{s}_{20.21^*}={}&\delta\log\mathrm{s}_{19.20}+\operatorname{cotang}u'\,.\,\delta u'-\operatorname{cotang}u''\,.\,\delta u''\\
\delta\log\mathrm{s}_{21^*.1^*}={}&\delta\log\mathrm{s}_{20.21^*}+\operatorname{cotang}v'\,.\,\delta v'-\operatorname{cotang}v''\,.\,\delta v''.
\end{aligned}
\]
\par\noindent\hspace*{\parindent}Daraus ergibt sich
\[
\begin{aligned}
&\Sigma(z_\mu-z_\lambda)\delta\log\mathrm{s}_{\lambda.\mu}=(z_{1^*}-z_1)\delta\log\mathrm{s}_{1.21}\\
&+(z_{1^*}-z_{21})\bigl(\operatorname{cotang}a'\,.\,\delta a'-\operatorname{cotang}a''\,.\,\delta a''\\
&\qquad+\operatorname{cotang}b'\,.\,\delta b'-\operatorname{cotang}b''\,.\,\delta b''\bigr)\\
&+(z_{1^*}-z_3)(\operatorname{cotang}c'\,.\,\delta c'-\operatorname{cotang}c''\,.\,\delta c'')+\ .\ .\ .\\
&+(z_{1^*}-z_{20})(\operatorname{cotang}u'\,.\,\delta u'-\operatorname{cotang}u''\,.\,\delta u'')\\
&+(z_{1^*}-z_{21^*})(\operatorname{cotang}v'\,.\,\delta v'-\operatorname{cotang}v''\,.\,\delta v''),
\end{aligned}
\]
\noindent oder, wenn man wie zuvor \(z_{1^*}\) und \(z_{21^*}\) wieder durch \(z_1\) und \(z_{21}\) ersetzt,
\[
\begin{aligned}
&3)\quad\Sigma(z_\mu-z_\lambda)\delta\log\mathrm{s}_{\lambda.\mu}=\\
&\quad(z_1-z_{21})\bigl(\operatorname{cotang}a'\,.\,\delta a'-\operatorname{cotang}a''\,.\,\delta a''\\
&\qquad+\operatorname{cotang}b'\,.\,\delta b'-\operatorname{cotang}b''\,.\,\delta b''\\
&\qquad+\operatorname{cotang}v'\,.\,\delta v'-\operatorname{cotang}v''\,.\,\delta v''\bigr)\\
&+(z_1-z_3)(\operatorname{cotang}c'\,.\,\delta c'-\operatorname{cotang}c''\,.\,\delta c'')\\
&+(z_1-z_4)(\operatorname{cotang}d'\,.\,\delta d'-\operatorname{cotang}d''\,.\,\delta d'')+\ .\ .\ .\\
&+(z_1-z_{20})(\operatorname{cotang}u'\,.\,\delta u'-\operatorname{cotang}u''\,.\,\delta u'').
\end{aligned}
\]
\par\noindent\hspace*{\parindent}Sind die Winkelkorrekturen in Bogensekunden angegeben, so folgt aus Gleichung (1)
\[
4)\quad\GaussOldenburgLeftBrace{\begin{aligned}
&206265\,.\,\delta z_{1^*}=\\
&\quad(z_1-z_{21})\bigl\{\operatorname{cotang}v'\,.\,\delta v'-\operatorname{cotang}v''\,.\,\delta v''\\
&\qquad+\operatorname{cotang}a'\,.\,\delta a'-\operatorname{cotang}a''\,.\,\delta a''\\
&\qquad+\operatorname{cotang}b'\,.\,\delta b'-\operatorname{cotang}b''\,.\,\delta b''\\
&\qquad+\mathrm{i}(\delta v^0+\delta a^0+\delta b^0)\bigr\}\\
&+(z_1-z_3)\{\operatorname{cotang}c'\,.\,\delta c'-\operatorname{cotang}c''\,.\,\delta c''-\mathrm{i}\,.\,\delta c^0\}\\
&+(z_1-z_4)\{\operatorname{cotang}d'\,.\,\delta d'-\operatorname{cotang}d''\,.\,\delta d''+\mathrm{i}\,.\,\delta d^0\}\\
&.\quad.\quad.\quad.\quad.\quad.\quad.\quad.\quad.\quad.\quad.\quad.\quad.\quad.\quad.\quad.\quad.\quad.\\
&+(z_1-z_{20})\{\operatorname{cotang}u'\,.\,\delta u'-\operatorname{cotang}u''\,.\,\delta u''-\mathrm{i}\,.\,\delta u^0\}.
\end{aligned}}
\]
\par\noindent\hspace*{\parindent}Damit die Punkte 1 und \(1^*\) zusammenfallen, muss man
\[
z_{1^*}+\delta z_{1^*}-z_1=0,\qquad \delta z_{1^*}=z_1-z_{1^*}=-\delta z_1
\]
\noindent setzen, wenn das vorangestellte \(\delta\) eine mit positivem Vorzeichen anzubringende Korrektur bezeichnet.
\par\endgroup
\clearpage
\noindent342\hfill{\small\textbf{BEMERKUNGEN.\quad DREIECKSKRANZ UM OLDENBURG.}}\par\medskip

\par\noindent\hspace*{\parindent}Die Seitengleichung 1) auf S.~336 ist die Bedingung dafür, dass sich bei der Rechnung durch das Kranzsystem der Wert der Seite \(\mathrm{s}_{1.21}\) wieder ergibt. Es gilt
\[
\frac{\mathrm{s}_{1^*.21^*}}{\mathrm{s}_{1.21}}=\frac{\sin a'}{\sin a''}\,.\,\frac{\sin b'}{\sin b''}\ .\ .\ .\frac{\sin v'}{\sin v''},
\]
\noindent also
\[
\begin{aligned}
&\delta\log\mathrm{s}_{1^*.21^*}=-\delta\log\mathrm{s}_{1.21}=-\frac{\delta z_1-\delta z_{21}}{z_1-z_{21}}\\
&\quad=\operatorname{cotang}a'\,.\,\delta a'-\operatorname{cotang}a''\,.\,\delta a''+\ .\ .\ .\\
&\quad+\operatorname{cotang}v'\,.\,\delta v'-\operatorname{cotang}v''\,.\,\delta v''.
\end{aligned}
\]
\par\noindent\hspace*{\parindent}Verwendet man den briggschen Logarithmus und sollen die Korrekturen \(\delta a',\ \delta a''\) usw. in Bogensekunden erhalten werden, so ist die linke Seite dieser Gleichung mit \(\frac{206265}{\mathrm{Mod.}}\) zu multiplizieren.

\par\noindent\hspace*{\parindent}In den Formeln 1) und 2) auf S.~336 und 337 steht im Original
\[
\frac{\delta(1-21)}{1-21}\ \text{und}\ \delta1\qquad\text{anstelle von}\quad-\frac{\delta(1-21)}{1-21}\ \text{und}\ -\delta1.
\]
\par\noindent\hspace*{\parindent}Im Jahr 1834 teilte \textsc{Gauss} (nach dem »Generalbericht über die mitteleuropäische Gradmessung für das Jahr 1865«, S.~22/23) dem oldenburgischen Vermessungsdirektor v.~\textsc{Schrenck} die 21 sphäroidischen, auf \(180^0+\text{Exzess}\) abgestimmten Dreiecke des Kranzes um Oldenburg mit. In einem Schreiben vom 3.~August 1835 äußert sich \textsc{Gauss} dazu wie folgt (nach dem genannten Generalbericht, S.~23/24).

\par\noindent\hspace*{\parindent}»Was nun die relative Genauigkeit der Dreiecke betrifft, so irren Euer Hochwohlgeboren, wenn Sie die Dreiecke 2--8 [in der Abbildung auf S.~332 sind dies die Dreiecke \(q,\ p,\ o,\ n,\ m,\ l,\ k\)] den übrigen 9--21 und 1 [in der Abbildung \(i\) bis \(a\), \(v\) bis \(r\)] gegenüberstellen. Der richtige Vergleich lautet, dass 3--16 [in der Abbildung \(p\) bis \(b\)] \textls[120]{viel} ungenauer sind als 17--21 [\(a,\ v,\ u,\ t,\ s\)], 1 [\(r\)] und 2 [\(q\)]. Die letzten 7 Dreiecke habe ich selbst mit 12-zölligen Theodoliten und größter Sorgfalt gemessen. Dabei diente ausnahmslos Heliotroplicht als Ziel, und ich achtete so weit wie möglich auf die Festigkeit der Standpunkte, von denen 3 zu ebener Erde lagen. Die 14 anderen Dreiecke wurden dagegen zu einem anderen Zweck, mit einem schwächeren Instrument (einem 8-zölligen Theodoliten), mit viel geringerem Zeitaufwand, ohne Heliotroplicht und mitunter an sehr ungünstigen Standpunkten gemessen, etwa auf den Türmen von Twistringen und Asendorf und vielleicht auch auf einigen der anderen Türme. Da ich also die ersten 7 Dreiecke für so genau gemessen halte, wie es der Stand der Kunst erlaubt, würde ich für die 14 übrigen nur \(\frac{1}{3}\) dieser Genauigkeit oder nur \(\frac{1}{9}\) dieses Gewichts ansetzen.«

\par\noindent\hspace*{\parindent}Bildet man die Winkelsumme des Siebenecks 4. 6. 10. 12. 16. 19. 21 aus den Winkeln, die \textsc{Gauss} v.~\textsc{Schrenck} mitteilte, so beträgt der Schlussfehler nach dem Generalbericht (S.~24) \(14\overset{\prime\prime}{,}900\).
\begin{flushright}\textsc{Krüger.}\end{flushright}
\begin{center}\rule{.26\textwidth}{.4pt}\end{center}
